\labsection{2}{Measuring financial inclusion}{sec:lab02-inclusion}

\begin{labproject}{Turn an inclusion claim into a number you can defend}
\textbf{Coverage.} Chapter~\ref{ch:financial-inclusion}.\\
\textbf{Goal.} Move from ``FinTech improves financial inclusion'' --- which is untestable as written --- to a specific, measured, falsifiable statement.

\courselabflow{Load sample}{Measure access}{Measure activity}{Inspect gaps}



\end{labproject}

\textbf{Complete worked example.} This runnable fixture uses synthetic inputs. First reproduce it, then adapt the data or ratings with documented evidence. All code below belongs to one file.
\begin{lstlisting}[style=utpterminal]
python all_experiments/lab02_measuring_inclusion.py
\end{lstlisting}

\noindent\textbf{Setup and output folder.}
\lstinputlisting[style=utppython,firstline=1,lastline=10]{all_experiments/lab02_measuring_inclusion.py}

\noindent\textbf{Generate a reproducible respondent fixture.}
\lstinputlisting[style=utppython,firstline=11,lastline=32]{all_experiments/lab02_measuring_inclusion.py}

\noindent\textbf{Measure access and subgroup denominators.}
\courselabtrace{Inclusion is a sequence of stricter sets}{Access}{An account exists}{Active}{Account exists and at least three recent transactions}{Meaningful}{Active account plus credit or savings usage}

\lstinputlisting[style=utppython,firstline=33,lastline=38]{all_experiments/lab02_measuring_inclusion.py}

\noindent\textbf{Define nested inclusion levels.}


\textbf{Use the chapter while coding.} Revisit \hyperref[ch:financial-inclusion]{the earlier chapter explanation}; match its operation to the variables and state changes below.

\textbf{Follow the supplied implementation.} Check one respondent through all three boolean expressions. The barrier analysis later changes the denominator to unbanked respondents. Its percentages cannot be compared as though they used the same population as the overall inclusion rates.
\courselabtrace{Trace the three implemented inclusion masks}{1. Input / state}{access = has\_account; active additionally requires at least three transactions in 90 days.}{2. Operation}{meaningful additionally requires credit or savings use, using OR between those two flags.}{3. Check}{The mean of each mask divides by all 600 synthetic respondents; rates must be nested.}


\lstinputlisting[style=utppython,firstline=39,lastline=47]{all_experiments/lab02_measuring_inclusion.py}

\noindent\textbf{Rank barriers within the appropriate population.}
\lstinputlisting[style=utppython,firstline=48,lastline=57]{all_experiments/lab02_measuring_inclusion.py}
\begin{labvisual}{Synthetic sample: progressively stricter inclusion}
\begin{labgraphic}\begin{tikzpicture}[font=\small]
\node[anchor=east,font=\scriptsize] at (-.2,0.0) {access};
\fill[UTPPaleBlue] (0,-0.16) rectangle (7,0.16);
\fill[UTPTeal] (0,-0.16) rectangle (4.5267,0.16);
\node[anchor=west,font=\scriptsize] at (7.15,0.0) {64.7\%};
\node[anchor=east,font=\scriptsize] at (-.2,-0.55) {active};
\fill[UTPPaleBlue] (0,-0.7100000000000001) rectangle (7,-0.39);
\fill[UTPTeal] (0,-0.7100000000000001) rectangle (2.5317,-0.39);
\node[anchor=west,font=\scriptsize] at (7.15,-0.55) {36.2\%};
\node[anchor=east,font=\scriptsize] at (-.2,-1.1) {meaningful};
\fill[UTPPaleBlue] (0,-1.26) rectangle (7,-0.9400000000000001);
\fill[UTPTeal] (0,-1.26) rectangle (1.4350,-0.9400000000000001);
\node[anchor=west,font=\scriptsize] at (7.15,-1.1) {20.5\%};
\end{tikzpicture}\end{labgraphic}\end{labvisual}

\textbf{Check the run.} The following output was captured from the supplied script; it describes the synthetic fixture only.
\lstinputlisting[style=utpterminal]{labs/outputs/lab02.txt}
% Reused from chapters/02-financial-inclusion.tex
\begin{labvisual}{Recall the chapter model}
\begin{labgraphic}
\begin{tikzpicture}[
  side/.style={draw=UTPBlue,fill=UTPPaleBlue,rounded corners=2mm,text width=5.4cm,minimum height=3.0cm,align=left,inner sep=3mm,font=\scriptsize\color{UTPNavy}},
  side2/.style={draw=UTPOrange,fill=UTPPaleOrange,rounded corners=2mm,text width=5.4cm,minimum height=3.0cm,align=left,inner sep=3mm,font=\scriptsize\color{UTPNavy}},
  center/.style={draw=UTPNavy,fill=UTPNavy,text=white,rounded corners=2mm,text width=2.4cm,minimum height=1.2cm,align=center,font=\small\bfseries},
  arr/.style={-{Stealth[length=2mm]},thick,draw=UTPMuted}
]
\node[side] (d) at (-4.0,0) {\textbf{Demand side}\par\smallskip
small/volatile income\par distance\par missing documents\par low literacy\par distrust\par negative history};
\node[center] (x) at (0,0) {FINANCIAL\\EXCLUSION};
\node[side2] (s) at (4.0,0) {\textbf{Supply side}\par\smallskip
high branch cost\par legacy IT\par manual processes\par standardized products\par weak competition\par limited innovation};
\draw[arr] (d.east)--(x.west); \draw[arr] (s.west)--(x.east);
\end{tikzpicture}
\end{labgraphic}\end{labvisual}


\begin{verification}
\textbf{Explain and extend your result.} Select two groups and report counts beside rates. Recalculate one measure with a different denominator and explain why the interpretation changes.

\begin{itemize}\item The script runs from a fresh process and creates the named CSV or report files.\item Every rate uses a stated population and denominator; every model assumption is identified.\item Synthetic outputs are labelled and are not presented as observations of a real institution.\item Change one input, rerun, and explain the direction of the result.\end{itemize}\end{verification}
\begin{labdeliverable}Submit the complete script, generated evidence, and a one-page interpretation linking the results to the chapter model. State the assumptions and one limitation.\end{labdeliverable}
